\begin{statement}{E.67}{Proved}{The original Mordell sign and the complete finite Gaussian norm comparison}
Retain the separate original orbifold actions $\Omega_3,\Omega_4$
and their Ferrari homology matrices $B_3,B_4$ from E.64.28--29.
For every multiplier-one normalized gentle pair, the action $B_3$
on the positive E.1 line $\mathbf F_3(1,-1)$ is the identity.
Consequently, for the full-GT comparison characters of E.64,
\[
 \epsilon_3(C_3)=\epsilon_3(D_3),\qquad
 \delta=\epsilon_3(D_3)+\epsilon_2(D_4)+\epsilon_2(B_4).
 \tag{E.67.1}
\]
The latter is an identity of the specified characters; it does
not assign the remaining sum on full GT.

The Gaussian term has a finite evaluation formula with its
actual word and origin retained. Put $V=\mathbf F_3^2$ and
\[
 R=\begin{pmatrix}0&1\\-1&1\end{pmatrix},\qquad
 S=\begin{pmatrix}0&1\\-1&0\end{pmatrix}.
 \tag{E.67.2}
\]
Use $(v,A)(w,B)=(v+Aw,AB)$ in the following two finite groups:
\[
 \begin{array}{c|c|c|c|c}
 &\text{group}&x&y&z\\ \hline
 \mathrm M&V\rtimes\langle R\rangle&(0,R)&(e_1,R^3)&(e_1,R^2)\\
 \mathrm G&V\rtimes\langle S\rangle&(e_1,S)&(-e_2,S^2)&(0,S).
 \end{array}
 \tag{E.67.3}
\]
These are the actual respective quotients of the two triangles
by $[K,K]K^3$, with orders $54$ and $36$. If $v$ is the
translation component of $f(x,y)$ in $\mathrm G$, then
\[
 \begin{array}{c|c|c}
 B_4\bmod3&v&\epsilon_2(B_4)\\ \hline
 I&(0,0)&0\\
 -I&(2,0)&0\\
 S&(1,1)&1\\
 -S&(1,2)&1
 \end{array}
 \qquad B_4e_1=e_1-v.
 \tag{E.67.4}
\]
In particular the bit is $v_2^2\in\{0,1\}\subset\mathbf F_3$,
read in $\mathbf F_2$.

There is also a complete finite correction-image comparison over
the actual E.1 source. Let $N=N_B$, $K_6=\overline{N'N^6}$,
and evaluate the same free word in the following ordered pairs,
first in $\mathrm M$ and then in $\mathrm G$:
\[
               (x,y),(y,x),(z,x),(x,z),(y,z),(z,y).
 \tag{E.67.5}
\]
Write $\Theta$ for this twelve-evaluation map. On $K_6$ its
values are translations. With $w=(1,-1)$ its exact image is
\[
 \begin{split}
 \Theta(K_6)=\Theta(N')&=\mathcal M\times\mathcal G,\\
 \mathcal M&=\{(aw,aw,bw,cw,-cw,-bw):a,b,c\in\mathbf F_3\},\\
 \mathcal G&=\{(u,v,t,-t,-Sv,-Su):u,v,t\in V\}.
 \end{split}
 \tag{E.67.6}
\]
Thus the common image has order $3^9$, and its two projection
ranks are $3$ and $6$. Every vector in this displayed product
has an actual word preimage in $N'\subset\widehat F'$.

In each of the $24$ original E.1 word fibres, and separately
in each of the $24$ E.17 word fibres used in E.62 and E.66,
there are exactly four attainable twelve-evaluation vectors
which satisfy all the two-cycle and three-cycle equations
and give determinant-one automorphisms of both groups in
\textup{(E.67.3)}. They have one uniquely specified framed
Mordell target and all four Gaussian targets in
\textup{(E.67.4)}. The two origin lists are evaluated separately;
no equality of their $G_6$ values is used.

These are complete correction-image fibres. For a retained
origin $f_j$ and one of its four values $t$, choose the saved
literal correction $c_{j,t}\in N'$. The entire corresponding
word fibre in $f_j(K_6\cap\widehat F')$ is exactly
\[
 f_jc_{j,t}\bigl(K_6\cap\widehat F'\cap\ker\Theta\bigr).
 \tag{E.67.7}
\]
The corrected pair induces an automorphism of the actual
two-factor image $\mathrm M\times\mathrm G$ of order $1944$,
so its map onto that image is generating. These assertions
retain their specified finite source and equations; they
do not assert a full-GT lift or an automorphism of an
unevaluated larger joint source.
\end{statement}
\begin{proof}
We prove the two quotient maps and their frame equations before
identifying the complete common correction image.

\paragraph{The actual Mordell quotient and its positive sign.}
The original triangle in E.64 has $xyz=1$, $y^2=z^3=1$.
Its free kernel $K$ has the Ferrari generators
\[
                 k_1=zx^{-2},\qquad k_2=x^4z.
 \tag{E.67.8}
\]
The cusp identity is $[k_1,k_2]x^6=1$. Hence $x^6=1$
after abelianizing $K$, and the deck action splits there.
Substitution in \textup{(E.67.3)} sends $k_1,k_2$ to
$e_1,e_2$, respectively. It sends the deck generator to
$R$, of order six. The resulting epimorphism has exactly
the kernel $[K,K]K^3$: the induced map on $K$ is its
abelianization modulo three and the deck quotient is unchanged.
This proves the stated order-$54$ comparison.

At multiplier one, E.64.28 fixes $x$ and sends
$z$ to $f(z,x)zf(z,x)^{-1}$. The conjugator has trivial
deck coordinate because $f\in\widehat F'$. Write its
translation as $a\in V$. The induced matrix $B_3$ commutes
with $R$, since $x$ is fixed. Comparing the two images of
$z$ gives the exact equality
\[
                    (B_3-I)e_1=(I-R^2)a.
 \tag{E.67.9}
\]
The image of $I-R^2$ is $\mathbf F_3w$, where
$w=(1,-1)$, and $e_1$ has nonzero image in $V/\mathbf F_3w$.
Direct solution of $BR=RB$ gives $B=A I+B_0R$.
Its scalar both on $\mathbf F_3w$ and on the indicated
one-dimensional quotient is $A-B_0$. Equality
\textup{(E.67.9)} makes that scalar one. This proves the
positive sign with the actual deck-zero frame retained.
The determinant-one condition from the primitive boundary
commutator in E.64.29 leaves the three matrices in this
positive centralizer coset.

The original Tate action used here is the separate map in
\cite[Proposition 4.4]{DE_NakamuraTsunogaiYasuda2010};
the comparison with $C_3$ is the actual quotient
$D_3=C_3B_3^{-1}$ of E.64. Thus its sign gives the first
identity in \textup{(E.67.1)} without identifying the two
actions. Applying the same identity $C_2=D_4B_4$ and then
the exact parity definition in E.65 gives the second.

\paragraph{The actual Gaussian quotient and all four finite actions.}
In the quartic subgroup of E.64 write $z=d$.
Its deck generator acts by $S$, and the retained Ferrari
expressions give
\[
 x=\xi_1d^{-3},\qquad y=\xi_2^{-1}d^2,
             \qquad d^4=1
 \tag{E.67.10}
\]
after taking the stated centre quotient. Abelianizing
$K=\langle\xi_1,\xi_2\rangle$ modulo three therefore
gives exactly the Gaussian row of \textup{(E.67.3)}.
In that quotient $x^2y=e_1$ and $z^2y=e_2$; these also
prove surjectivity and the exact kernel comparison.

For an actual multiplier-one orbifold action, $B_4$ commutes
with $S$ and has determinant one by E.64.29. Solving these
two equations over $\mathbf F_3$ gives exactly
$I,-I,S,-S$. Put $f(x,y)=(v,I)$. Its two-cycle partner
is $f(y,x)=(-v,I)$, so the action sends $y$ to
$(-v,I)y(v,I)$. Evaluating $x^2y=e_1$, with $x$ fixed,
gives $B_4e_1=e_1-v$. This proves the complete table
\textup{(E.67.4)}, including the parity and its sign.

All four finite automorphisms have explicit frames. Given
$B\in\{I,-I,S,-S\}$, fix $x$ and put
\[
        z_B=((I-B)e_1,S),\qquad y_B=(z_Bx)^{-1}.
 \tag{E.67.11}
\]
The resulting map acts by $B$ on $V$. Its multiplication
law is preserved, and its inverse uses $B^{-1}$, so it is
an automorphism of the full order-$36$ group. Equivalently,
the certificate checks its entire permutation on those
$36$ elements. The matrices $I-S$ and $I-S^2$ are
invertible, so the conjugating translations $h_z,h_y$
with $z_B=h_zzh_z^{-1}$ and $y_B=h_yyh_y^{-1}$ are
unique.

In the Mordell group, a positive determinant-one matrix
$B$ instead gives $y_B=(Be_1,R^3)$ and
$z_B=(Be_1,R^2)$. The $y$ frame is unique because
$I-R^3$ is invertible. The $z$ frames form a coset of
$\ker(I-R^2)$, of order three. Consequently there are
exactly nine framed Mordell targets, over its three
possible matrices. For either group the complete six
translation values dictated by these frames and the
two-cycle and three-cycle equations are
\[
 \bigl(-h_y,\ h_y,\ h_z,\ -h_z,\ h_y-h_z,\ h_z-h_y\bigr).
 \tag{E.67.12}
\]
This supplies all $9\cdot4$ candidate pairs before
intersecting with any old-source fibre; the $z$-frame
ambiguity in the Mordell factor is fully retained.

\paragraph{The complete actual correction image.}
Use the $97$ literal free generators $n_i$ of $N_B$ from
the previously certified E.1 source, with the unchanged
word list in the present input. In every Mordell
substitution their images have even deck coordinate,
so lie in $V\rtimes\langle R^2\rangle$. This group
has exponent three: $I+R^2+R^4=0$. In every Gaussian
substitution their images lie in
$V\rtimes\langle S^2\rangle=V\rtimes\langle-I\rangle$,
whose exponent divides six. Thus every sixth power
from $N_B$ is killed in all twelve maps. This proves
$\Theta(K_6)=\Theta(N_B')$, including the closed
profinite subgroups.

Every commutator image is a translation. The complete
calculation evaluates all $\binom{97}{2}=4656$ literal
commutators, retains a word preimage for each pivot,
and closes their span under each of the $97$ actual
conjugation matrices. The initial span already has
dimension nine and is stable under all these conjugations.
It is therefore the entire image of $N_B'$, since those
commutators normally generate $N_B'$. The saved nine
pivot words are each replayed in all twelve maps.

The two projection ranks are three and six. More
explicitly, row reduction gives exactly
$\mathcal M$ and $\mathcal G$ in
\textup{(E.67.6)}. The separate small linear certificate
checks both inclusions between the nine saved rows and
the displayed parametrization. Since the common rank is
the sum of the projection ranks, the image is their
full product on this actual common domain. In particular
every Gaussian correction can be made while preserving
all six Mordell values.

There is an additional exact upper-bound check on the
Gaussian side. Its displayed six-dimensional space
contains the value of $[x,y]$ and is stable under the
two actual free-generator conjugations in the six
substitutions. It therefore contains the image of the
whole $\widehat F'$. The already computed image of
$N_B'$ is the entire space, proving equality. This
uses the literal free-source maps, rather than an
identification with one of the cyclic distribution
maps used elsewhere.

\paragraph{Complete finite fibres and generating maps.}
The homogeneous two-cycle and three-cycle equations
have zero intersection with $\mathcal M$. Indeed in
\textup{(E.67.6)} the first two-cycle equation gives
$a=0$, the second gives $c=-b$, and the three-cycle
equation gives $b-c=0$. Hence $a=b=c=0$. Thus any
given origin has at most one framed Mordell norm
target. Testing the complete nine-target list against
the actual affine origin coset gives one in each of
the $24$ E.1 cases and, independently, in each of the
$24$ E.17 cases.

For the Gaussian space the homogeneous norm equations
give $v=-u$ and $t=-(I+S)u$. Intersecting with the
determinant-one automorphism condition gives exactly
the four values of $u=f(x,y)$ in
\textup{(E.67.4)}. Each of the actual origin vectors
lies in $\mathcal G$, so all four absolute targets
are reached. Together with the full-product equality
this proves the four-point correction-image fibre.
The certificate also constructs and replays one
literal $N_B'$ correction for every one of the
resulting $96$ values in each origin list. Applying
the homomorphism $\Theta$ to two such words proves
the exact coset formula \textup{(E.67.7)}.

The literal pair $(x_{\mathrm M},x_{\mathrm G})$,
$(y_{\mathrm M},y_{\mathrm G})$ generates the full
product $\mathrm M\times\mathrm G$: enumeration of
its actual subgroup gives $54\cdot36=1944$ elements.
For every retained corrected word, the new two-factor
pair is the image of this pair under the product of
the two constructed automorphisms. Thus the corrected
source map is also onto this specified product.

\paragraph{Preserved attempts and reusable maps.}
Some original E.1 representatives fail a Mordell
three-cycle equation before correction; their indices
and values are retained. Other fixed representatives
fail the Gaussian action or its determinant-one
condition. The proof above uses the complete affine
correction image, so none of those fixed-word failures
is used to exclude its full fibre.

The authoritative programs share the prefix
\texttt{bc\_orbifold\_CM\_}: their suffixes are
\texttt{characters.py}, \texttt{six.py},
\texttt{linear.py}, and \texttt{export.py}.
The export retains the two literal free-generator
permutations on $108$ points, all $97$ original
$N_B$ permutations, the nine correction pivots,
the complete framed targets, and both distinct
origin lists. Its right-action convention is
$p(gh)=p(h)[p(g)]$; the affine action is
$t\mapsto A^{-1}(t-v)$. Each literal source word
is replayed in that convention. This gives exact
inputs for further common-source intersections,
without assigning a common quotient with another
three-primary cover.
\end{proof}
